% =====================================================================
%  CARNET D'ENTRAÎNEMENT — FICHE 5 (Bilans de matière)
%  THÈME 5 — Rendement et bilan de matière complet
%  Fiche TUTO
% =====================================================================
\documentclass[11pt,a4paper]{article}
\usepackage[utf8]{inputenc}
\usepackage[T1]{fontenc}
\usepackage{lmodern}
\usepackage[french]{babel}
\usepackage{amsmath,amssymb,mathtools}
\usepackage{icomma}
\usepackage[version=4]{mhchem}
\usepackage{siunitx}
\sisetup{output-decimal-marker={,},per-mode=symbol}
\DeclareSIUnit\Molar{\mole\per\litre}
\usepackage{xcolor}
\usepackage{colortbl}
\usepackage{tikz}
\usepackage[most]{tcolorbox}
\usepackage{enumitem}
\usepackage{array}
\usepackage{booktabs}
\usepackage{fancyhdr}
\usepackage[a4paper,top=1.5cm,bottom=1.4cm,left=1.8cm,right=1.8cm,headheight=24pt,headsep=6pt]{geometry}
\usetikzlibrary{arrows.meta,positioning,calc}

\definecolor{primary}{RGB}{142,93,168}
\definecolor{primarydark}{RGB}{82,45,108}
\definecolor{defcol}{RGB}{0,114,178}
\definecolor{thmcol}{RGB}{0,140,120}
\definecolor{formulacol}{RGB}{198,93,9}
\definecolor{attncol}{RGB}{200,40,55}
\definecolor{excol}{RGB}{0,135,90}
\newcommand{\Vm}{V_{\mathrm m}}

\pagestyle{fancy}\fancyhf{}
\fancyhead[L]{\small\textcolor{primarydark}{\textbf{Fiche 5} — Bilans de matière}}
\fancyhead[R]{\small\textcolor{primarydark}{\textsc{Tuto — Thème 5}}}
\fancyfoot[C]{\small\thepage}
\renewcommand{\headrulewidth}{0.8pt}
\renewcommand{\headrule}{\hbox to\headwidth{\color{primary}\leaders\hrule height \headrulewidth\hfill}}

\newtcolorbox{tutobox}[2][]{enhanced,breakable,boxrule=0.9pt,arc=2.5pt,
  left=3mm,right=3mm,top=2mm,bottom=2mm,colback=#2!6,colframe=#2,
  fonttitle=\bfseries,coltitle=white,title={#1}}

\setlength{\parindent}{0pt}
\setlength{\parskip}{2.3pt}

\begin{document}

\begin{center}
{\LARGE\bfseries\color{primarydark}Thème 5 — Rendement \& bilan complet}\\[1pt]
{\color{primary}\rule{0.5\linewidth}{1.2pt}}\\[2pt]
{\itshape La synthèse de toute la fiche : de la réalité expérimentale (rendement) au dosage.}
\end{center}
\vspace{-3pt}

\begin{tutobox}[Rendement théorique et rendement réel]{defcol}
Jusqu'ici on supposait que \emph{tout} le réactif limitant se transformait (rendement \SI{100}{\percent}).
En réalité, des réactions secondaires, un équilibre non total ou des pertes font qu'on obtient
\textbf{moins} de produit que prévu. Le \textbf{rendement} $R$ compare la quantité \emph{réellement}
obtenue à la quantité \emph{théorique} (maximale, fixée par le limitant).
\end{tutobox}

\begin{tutobox}[La formule du rendement]{formulacol}
\[
\boxed{\;R=\dfrac{n_{\text{réel}}}{n_{\text{théorique}}}\times 100\;}\quad(\si{\percent})
\qquad\Longleftrightarrow\qquad
n_{\text{réel}}=\dfrac{R}{100}\times n_{\text{théorique}}
\]
\textbf{où} $n_{\text{théorique}}$ est calculé par la stœchiométrie \emph{à partir du réactif limitant},
et $n_{\text{réel}}$ est \emph{mesuré}. On peut aussi l'écrire avec des masses
($R=m_{\text{réel}}/m_{\text{théorique}}$), car $m=nM$ et $M$ se simplifie. $R$ est compris entre 0 et \SI{100}{\percent}.
\end{tutobox}

\begin{tutobox}[L'organigramme universel d'un bilan de matière]{primary}
\begin{enumerate}[leftmargin=1.6em,topsep=0pt,itemsep=0.5pt]
  \item \textbf{Pondérer} l'équation (Thème 1) — aucun calcul valable sinon.
  \item \textbf{Convertir} toutes les données en moles (Thème 2 : $n=m/M$, $n=V/\Vm$, $n=CV$).
  \item \textbf{Identifier le limitant} : plus petit $n_i/\nu_i$ (Thème 4) ; il fixe $\xi_{\max}$.
  \item \textbf{Calculer les quantités finales} (produits, excès) par les coefficients (Thème 3).
  \item \textbf{Appliquer le rendement} si $R\neq\SI{100}{\percent}$ : $n_{\text{réel}}=\frac{R}{100}n_{\text{théo}}$.
  \item \textbf{Reconvertir} dans l'unité demandée ($m=nM$, $V=n\Vm$, $C=n/V$) et conclure.
\end{enumerate}
\end{tutobox}

\begin{tutobox}[Le titrage (dosage à l'équivalence)]{thmcol}
Doser, c'est déterminer une quantité \emph{inconnue} en la faisant réagir avec une quantité
\emph{connue} de réactif titrant. \textbf{À l'équivalence}, les réactifs ont été introduits dans les
\emph{proportions stœchiométriques exactes} : ni excès, ni défaut. Si le \emph{titré} a le coefficient
$a$ et le \emph{titrant} le coefficient $b$ :
\[
\boxed{\;\dfrac{n_{\text{titré}}}{a}=\dfrac{n_{\text{titrant}}}{b}\;}\qquad\text{avec } n=C\times V .
\]
On en tire la concentration inconnue : $C_{\text{titré}}=\dfrac{n_{\text{titré}}}{V_{\text{titré}}}$.
\end{tutobox}

\begin{tutobox}[Les pièges à éviter (rappel)]{attncol}
\begin{itemize}[leftmargin=1.3em,topsep=0pt,itemsep=0pt]
  \item Un \textbf{excès} de réactif \emph{ne garantit pas} \SI{100}{\percent} de rendement : il permet
        seulement de consommer tout le limitant, pas d'annuler les pertes.
  \item Un résultat \textbf{en moles} n'est ni une masse ($\times M$) ni un volume ($\times\Vm$).
  \item \textbf{mL $\to$ L}, \textbf{mg $\to$ g} : convertir avant tout calcul.
\end{itemize}
\end{tutobox}

\end{document}
